% ==============================================================================
% CAPÍTULO 1 - INTRODUÇÃO
% ==============================================================================

\section{Introdução}
\label{sec:introducao}

Agentes de código baseados em grandes modelos de linguagem (LLMs, do inglês
\textit{Large Language Models}) deixaram de completar linhas isoladas para
operar diretamente sobre repositórios: leem arquivos, editam código, executam
testes e iteram até considerar uma tarefa concluída. Ferramentas como o Claude
Code \cite{anthropic2025claudecode} e o OpenCode \cite{opencode2025docs}
materializam esse modo de trabalho em linha de comando, e \textit{benchmarks}
como o SWE-bench \cite{jimenez2024swebench} mostram que a taxa de resolução de
tarefas reais cresce com a capacidade do modelo. Essa capacidade, porém, tem
preço: os modelos de ponta custam mais por \textit{token}, dependem de cotas de
uso e nem sempre podem ser empregados, seja por restrições de custo, de
privacidade dos dados ou pela necessidade de execução local em
\textit{hardware} modesto.

Em paralelo, consolidou-se a prática de os próprios repositórios carregarem uma
estrutura de apoio para os agentes que neles trabalham. Arquivos de instrução
como \texttt{AGENTS.md} \cite{agentsmd2025} e \texttt{CLAUDE.md}
\cite{anthropic2025bestpractices} descrevem convenções, comandos e regras do
projeto; \textit{skills} e modelos de código ficam disponíveis para leitura sob
demanda; e arquivos de configuração habilitam servidores de linguagem (LSP,
\textit{Language Server Protocol}) e outras ferramentas. Neste trabalho, esse
conjunto de artefatos versionados no repositório é chamado de
\textit{harness}\footnote{O termo \textit{harness}, sem tradução consolidada em
português, designa aqui o conjunto de regras, \textit{skills} e configuração de
ferramentas que o repositório oferece ao agente. A expressão ``andaimes''
(\textit{scaffolding}) do título refere-se ao mesmo conjunto, visto como
estrutura de apoio ao modelo.}. A hipótese implícita nessa prática, resumida no
lema \textit{scaffolding over scale}, é que uma boa estrutura de apoio compensa,
ao menos em parte, a escala do modelo: um modelo menor bem orientado produziria
resultados comparáveis aos de um modelo maior sem orientação. Apesar da adoção
ampla, a evidência controlada sobre esse efeito ainda é escassa
\pendente{CITAÇÃO: estudos empíricos sobre o efeito de arquivos de instrução de
repositório (AGENTS.md, CLAUDE.md, regras do Cursor) em agentes de código}.

O problema que orienta esta pesquisa é, portanto, o seguinte: um
\textit{harness} em nível de repositório compensa a escala do modelo? A
pergunta desdobra-se em três: (i) o \textit{harness} aumenta a taxa com que
agentes resolvem tarefas reais, medida por testes que o agente não vê?
(ii) o \textit{harness} muda a qualidade do código produzido, mesmo quando a
tarefa já é resolvida? e (iii) qual camada do \textit{harness}, se alguma,
carrega o efeito observado?

O objetivo geral deste trabalho é avaliar, por meio de um experimento de
ablação, o efeito de um \textit{harness} de repositório real sobre a saída de
agentes de código de diferentes escalas, em tarefas derivadas do histórico
desse mesmo repositório. Os objetivos específicos são:

\begin{enumerate}[nosep]
    \item construir um \textit{benchmark} que reproduza \textit{pull requests}
    reais de uma aplicação Laravel, usando os testes de cada PR como oráculo
    oculto;
    \item definir níveis cumulativos de \textit{harness} a partir de um
    \textit{snapshot} congelado do repositório e aplicá-los de forma idêntica a
    todas as tarefas;
    \item executar seis modelos de escalas distintas em todas as tarefas e
    níveis, em ambiente isolado e reprodutível, registrando acerto, custo e
    comportamento;
    \item medir a qualidade do código produzido em três camadas independentes:
    análise estática, aderência às convenções do repositório e revisão por
    juízes LLM às cegas;
    \item analisar estatisticamente o efeito de cada nível, pareado por tarefa
    dentro de cada modelo, e discutir os resultados à luz da hipótese
    \textit{scaffolding over scale}.
\end{enumerate}

As contribuições do trabalho são três. A primeira é um \textit{benchmark} de
ablação sobre PRs reais com oráculo oculto, cujo repositório-sujeito é privado
e, portanto, pouco sujeito à contaminação dos dados de treinamento dos modelos.
A segunda é um orquestrador reprodutível que executa a matriz
modelo~$\times$~tarefa~$\times$~nível em \textit{sandbox} Docker, com
isolamento da árvore de trabalho, detecção de contaminação, tratamento de
limites de taxa e retomada sem duplicar células. A terceira é uma análise de
qualidade de código em três camadas, com juízes LLM de famílias distintas da do
modelo avaliado, validação cruzada entre camadas e medidas de concordância
entre juízes. Os dados resultantes compreendem 309 execuções avaliadas, de um
desenho de 336 células, e 615 revisões de juízes.

Antecipando os resultados, o experimento mostrou que quatro dos seis modelos
resolveram 95\% ou mais das tarefas já sem nenhum \textit{harness}, que nenhum
nível alterou a taxa de acerto de forma detectável e que o \textit{harness}
encareceu cada execução. Em contrapartida, o \textit{harness} elevou de forma
consistente as notas de convenções e de testes atribuídas pelos juízes e o
número de casos de teste escritos, sem alterar requisitos, segurança ou escopo.
O efeito concentrou-se nos modelos de desempenho intermediário e não apareceu
nos mais fracos.

O restante do artigo está organizado da seguinte forma. A
\autoref{sec:metodologia} classifica a pesquisa e descreve os procedimentos, as
métricas, os testes estatísticos e a mudança de planejamento ocorrida durante o
experimento. A \autoref{sec:referencial-teorico} situa o trabalho em cinco
eixos: \textit{benchmarks} de agentes de código, engenharia de contexto,
ferramentas de agentes, LLM como avaliador e métricas de qualidade de código. A
\autoref{sec:desenvolvimento} apresenta o objeto de estudo, o \textit{benchmark}
e o orquestrador, e em seguida os resultados de \textit{pass@1}, custo e
qualidade, a discussão e as ameaças à validade. A \autoref{sec:consideracoes}
retoma os objetivos, resume as contribuições e as limitações e aponta trabalhos
futuros.
